%%%%%%%%%%%%%%%%%%%%%%%%%%%%%%%%%%%%%%%%%%%%%%%%%%%%%%%%%%%%%%%%%%%%%%%%
% Plantilla TFG/TFM
% Escuela Politécnica Superior de la Universidad de Alicante
% Realizado por: Jose Manuel Requena Plens
% Contacto: info@jmrplens.com / Telegram:@jmrplens
%%%%%%%%%%%%%%%%%%%%%%%%%%%%%%%%%%%%%%%%%%%%%%%%%%%%%%%%%%%%%%%%%%%%%%%%

\chapter*{Preámbulo}
\thispagestyle{empty}
\begin{quote}
Este trabajo surge del interés por explorar los límites del razonamiento lingüístico en los modelos de lenguaje actuales. El objetivo es diseñar y evaluar estrategias de prompting que permitan a estos modelos resolver puzzles de olimpiadas lingüísticas en idiomas nunca vistos durante su entrenamiento. El trabajo se enmarca en el área de procesamiento de lenguas de baja densidad de recursos y ha sido realizado bajo la supervisión de Juan Antonio Pérez Ortiz en el Grado de Ingeniería Informática de la Universidad de Alicante.
\end{quote}


\cleardoublepage %salta a nueva página impar
\chapter*{Agradecimientos
}

\thispagestyle{empty}
\vspace{1cm}

Quisiera agradecer, en primer lugar, a mi tutor, Juan Antonio Pérez Ortiz, por su orientación, sus consejos y su disponibilidad durante el desarrollo de este Trabajo de Fin de Grado. Su ayuda ha sido fundamental para dar forma al proyecto y para afrontar las distintas decisiones que han ido surgiendo a lo largo del proceso.

También quiero dar las gracias a mi familia, por su apoyo constante durante todos estos años. Su confianza, paciencia y ánimo han sido esenciales no solo durante la realización de este trabajo, sino a lo largo de toda mi etapa universitaria.

Este trabajo nace también de un interés personal por los idiomas y por la forma en que las lenguas permiten entender distintas culturas. Haber vivido en diferentes países me ha permitido acercarme a varios idiomas desde una perspectiva práctica y humana, y ha despertado en mí una curiosidad especial por la traducción y el razonamiento entre lenguas. Esa experiencia personal ha influido directamente en la elección de este tema, centrado en la resolución de problemas de olimpiadas lingüísticas mediante modelos de inteligencia artificial.

Por último, quiero agradecer a todas las personas que, de una forma u otra, me han acompañado durante este camino: compañeros, amigos y profesores que han contribuido a mi formación y me han ayudado a llegar hasta aquí.

%\cleardoublepage %salta a nueva página impar
% Aquí va la dedicatoria si la hubiese. Si no, comentar la(s) linea(s) siguientes
%\chapter*{}
%\setlength{\leftmargin}{0.5\textwidth}
%\setlength{\parsep}{0cm}
%\addtolength{\topsep}{0.5cm}
%\begin{flushright}
%\small\em{
%A mi esposa Marganit, y a mis hijos Ella Rose y Daniel Adams,\\
%sin los cuales habría podido acabar este libro dos años antes \footnote{Dedicatoria de Joseph J. Roman en "An Introduction to Algebraic Topology"}
%}
%\end{flushright}


\cleardoublepage %salta a nueva página impar
% Aquí va la cita célebre si la hubiese. Si no, comentar la(s) linea(s) siguientes
\chapter*{}
\setlength{\leftmargin}{0.5\textwidth}
\setlength{\parsep}{0cm}
\addtolength{\topsep}{0.5cm}
\begin{flushright}
\small\em{
El límite de mi lenguaje es el límite de mi mundo.
}
\end{flushright}
\begin{flushright}
\small{
Ludwig Wittgenstein.
}
\end{flushright}
\cleardoublepage %salta a nueva página impar
